% Generated draft: requires mathematical and source review.
\documentclass[11pt]{article}
\usepackage[margin=25mm]{geometry}
\usepackage{fontspec}
\setmainfont{Latin Modern Roman}
\usepackage{amsmath,amssymb,mathtools,unicode-math}
\setmathfont{Latin Modern Math}
\usepackage{enumitem,needspace,xurl,hyperref}
\setlength{\parindent}{0pt}
\setlength{\parskip}{5pt}
\setlength{\emergencystretch}{3em}
\newcommand{\N}{\mathbb{N}}
\newcommand{\Z}{\mathbb{Z}}
\newcommand{\Q}{\mathbb{Q}}
\newcommand{\R}{\mathbb{R}}
\newcommand{\C}{\mathbb{C}}
\newcommand{\F}{\mathbb{F}}
\newcommand{\A}{\mathbb{A}}
\newcommand{\PP}{\mathbb{P}}
\newcommand{\E}{\mathbb{E}}
\newcommand{\eps}{\varepsilon}
\newcommand{\dd}{\,\mathrm{d}}
\begin{document}
\section*{UnsolvedMath ID: 701380}
\textbf{Source problem code:} catalogue entry
\subsection{Definitions and mathematical statement}
Let \(\mathfrak G\) denote the variety of all groups. For a group \(G\), let \(\operatorname{var}(G)\) be the smallest group variety containing \(G\), equivalently the class of groups satisfying every group identity satisfied by \(G\). This variety is proper if \(\operatorname{var}(G)\ne\mathfrak G\).

Under the standard interpretation of ``without free subsemigroups,'' require that \(G\) contain no subsemigroup isomorphic to the free semigroup on two generators. Explicitly, for every \(a,b\in G\), there are distinct nonempty words \(u,v\) in two formal letters such that \(u(a,b)=v(a,b)\).

The supplied record does not define \(\mathfrak A_p\mathfrak A\) or specify or quantify \(p\). Under the conventional interpretation for a prime \(p\), let \(\mathfrak A\) be the variety of abelian groups and \(\mathfrak A_p\) the variety of abelian groups of exponent dividing \(p\). Their product \(\mathfrak A_p\mathfrak A\) consists of groups \(H\) having a normal subgroup \(N\trianglelefteq H\) with \(N\in\mathfrak A_p\) and \(H/N\in\mathfrak A\).

With these conventions and with \(p\) treated as a fixed prime parameter, the question is whether there exists a finitely generated group \(G\) such that
\[
G\text{ contains no free subsemigroup on two generators},
\qquad
\mathfrak A_p\mathfrak A\subseteq\operatorname{var}(G)\subsetneq\mathfrak G.
\]
Here ``containing'' is read as non-strict inclusion. The intended range and quantification of \(p\), and the conventions for the undefined notation and free subsemigroups, require confirmation from the original problem. Thus this is a conditional formulation of the supplied question, not a resolution of those omissions.
\subsection{Short English statement}
Under the stated conventions, the question asks for a finitely generated group with no free semigroup on two generators whose generated variety is proper and contains every group with an abelian normal subgroup of exponent dividing \(p\) and abelian quotient. The source leaves \(p\) unspecified. It labels the problem open; that label has not been independently verified.
\subsection{Sources}
\begin{itemize}
\item Local UnsolvedMath registry, permanent ID \(701380\), exact source problem code \texttt{catalogue entry}, title ``Kourovka Notebook Problem 17.78.'' Supplied upstream revision: \texttt{d371ee917ac9647fbb06cd6fc529deec6f3dda82}.
\item Supplied attribution: O. Macedo\'nska, proposed in 2010. The supplied background identifies \emph{The Kourovka Notebook (Unsolved Problems in Group Theory)}, 17th issue (2010), original problem 17.78, and the 21st issue edition (2026), arXiv:1401.0300v48. The collection description names E. I. Khukhro and V. D. Mazurov as editors.
\item The local record was created and updated on 2026-10-16 and carries the status ``open.'' These are supplied registry metadata, not independent source verification or a verified literature assessment. The cited Notebook text was not independently consulted for this draft.
\end{itemize}
\end{document}
